% GENERATED FILE. Edit canonical lessons, then run npm run generate:books.
% canonical-source-hash: fnv1a64:f1079454af4c7448
% canonical-lessons: BN-C308-bisorjon, BN-C308-bhasano, BN-C308-biday, BN-C308-shobhajatra, BN-C308-kolakuli-kora

\chapter{The Last Day of the Festival}
\label{ch:bn-the-last-day-of-the-festival}
\hlchaptermodality{308}

\begin{chapteropening}
\textbf{By the end of this chapter:} \emph{I can tell how a festival ends: the immersion, the farewell, the embrace.}

Complete the last lesson of chapter 308: I can tell how a festival ends: the immersion, the farewell, the embrace.
\end{chapteropening}

\section[\texorpdfstring{\bn{বিসর্জন} (bisôrjôn)}{বিসর্জন (bisôrjôn)}]{\bn{বিসর্জন} (bisôrjôn) — the immersion of an idol}

\label{lesson:BN-C308-bisorjon}



\begin{quote}
Before the new one: say the Bengali for an evening worship with lamps, then the Bengali for to swing.
\end{quote}

\subsection*{You'll want to know: \bn{বিসর্জন}}
\textbf{\bn{বিসর্জন}} — \emph{bisôrjôn} — \textquotedblleft{}the immersion of an idol\textquotedblright{}.

\begin{quote}
\bn{পুজোর} \bn{শেষ} \bn{দিনে} \bn{প্রতিমার} \bn{বিসর্জন} \bn{হয়।}
\end{quote}

\begin{quote}
\emph{pujor shesh dine prôtimār bisôrjôn hôy.}
\end{quote}

\begin{quote}
On the last day of the puja comes the immersion of the image.
\end{quote}

\begin{cousinweb}
\textbf{\bn{বিসর্জন}} (\emph{bisôrjôn}) is Sanskrit \emph{visarjana}, 'sending forth, dismissal' (Monier-Williams), which already names the casting of a deity's image into holy water at the close of a festival.
\end{cousinweb}

\begin{grammarlens}[title={What you've built}]
One more word for telling what happened.
\end{grammarlens}

\subsection*{Your turn}
\begin{itemize}
\raggedright
  \item \emph{Say it:} \emph{bisôrjôn}
  \item \emph{Say it:} \emph{bisôrjôn}, once more
  \item \emph{Say it:} \emph{pujor shesh dine prôtimār bisôrjôn hôy.} — On the last day of the puja comes the immersion of the image.
  \item \emph{Recall:} say the Bengali for an evening worship with lamps, then the Bengali for to swing, then say \emph{bisôrjôn} again
\end{itemize}

\begin{tcolorbox}[breakable,colback=teal!4,colframe=teal!35!black,title={Before you move on}]
What does \bn{বিসর্জন} mean? (\textquotedblleft{}The immersion of an idol\textquotedblright{}.) Say it once more.
\end{tcolorbox}
\section[\texorpdfstring{\bn{ভাসানো} (bhāsāno)}{ভাসানো (bhāsāno)}]{\bn{ভাসানো} (bhāsāno) — to set afloat; to immerse}

\label{lesson:BN-C308-bhasano}



\begin{quote}
Before the new one: say the Bengali for to swing, then the Bengali for the immersion of an idol.
\end{quote}

\subsection*{You'll want to know: \bn{ভাসানো}}
\textbf{\bn{ভাসানো}} — \emph{bhāsāno} — \textquotedblleft{}to set afloat; to immerse\textquotedblright{}.

\begin{quote}
\bn{সন্ধ্যায়} \bn{তারা} \bn{প্রতিমা} \bn{নদীতে} \bn{ভাসাল।}
\end{quote}

\begin{quote}
\emph{sôndhyāy tārā prôtimā nôdīte bhāsālô.}
\end{quote}

\begin{quote}
In the evening they set the image afloat on the river.
\end{quote}

Past (he/she …): \bn{সেই} \bn{শিশু} \bn{কাগজের} \bn{নৌকা} \bn{ভাসাল}, \emph{sei shishu kāgôjer noukā bhāsālô}.

\begin{cousinweb}
\textbf{\bn{ভাসানো}} (\emph{bhāsāno}) is \textbf{\bn{ভাসা}}, to float — CDIAL 9654, Sanskrit \emph{bhráśyati}, 'falls', which lists Bengali \emph{bhāsā}, 'to float, drift along' — with the -ānō ending: to set something afloat.
\end{cousinweb}

\begin{grammarlens}[title={What you've built}]
One more word for telling what happened.
\end{grammarlens}

\subsection*{Your turn}
\begin{itemize}
\raggedright
  \item \emph{Say it:} \emph{bhāsāno}
  \item \emph{Say it:} \emph{bhāsāno}, once more
  \item \emph{Say it:} \emph{sôndhyāy tārā prôtimā nôdīte bhāsālô.} — In the evening they set the image afloat on the river.
  \item \emph{Recall:} say the Bengali for to swing, then the Bengali for the immersion of an idol, then say \emph{bhāsāno} again
\end{itemize}

\begin{tcolorbox}[breakable,colback=teal!4,colframe=teal!35!black,title={Before you move on}]
What does \bn{ভাসানো} mean? (\textquotedblleft{}To set afloat; to immerse\textquotedblright{}.) Say it once more.
\end{tcolorbox}
\section[\texorpdfstring{\bn{বিদায়} (bidāy)}{বিদায় (bidāy)}]{\bn{বিদায়} (bidāy) — a farewell; a parting}

\label{lesson:BN-C308-biday}



\begin{quote}
Before the new one: say the Bengali for the immersion of an idol, then the Bengali for to set afloat.
\end{quote}

\subsection*{You'll want to know: \bn{বিদায়}}
\textbf{\bn{বিদায়}} — \emph{bidāy} — \textquotedblleft{}a farewell; a parting\textquotedblright{}.

\begin{quote}
\bn{বিদায়ের} \bn{সময়} \bn{দিদিমা} \bn{কাঁদলেন।}
\end{quote}

\begin{quote}
\emph{bidāyer sômôy didimā kā̃dlen.}
\end{quote}

\begin{quote}
Grandmother cried when it was time to say goodbye.
\end{quote}

\begin{cousinweb}
\textbf{\bn{বিদায়}} (\emph{bidāy}) is Sanskrit \emph{vidāya}; Monier-Williams gives 'dismissal with good wishes' and suspects that this sense is not Sanskrit at all but Arabic \emph{wadāʿ}, 'farewell'. \bn{বিদায়} \bn{নেওয়া}, to take your leave.
\end{cousinweb}

\begin{grammarlens}[title={What you've built}]
One more word for telling what happened.
\end{grammarlens}

\subsection*{Your turn}
\begin{itemize}
\raggedright
  \item \emph{Say it:} \emph{bidāy}
  \item \emph{Say it:} \emph{bidāy}, once more
  \item \emph{Say it:} \emph{bidāyer sômôy didimā kā̃dlen.} — Grandmother cried when it was time to say goodbye.
  \item \emph{Recall:} say the Bengali for the immersion of an idol, then the Bengali for to set afloat, then say \emph{bidāy} again
\end{itemize}

\begin{tcolorbox}[breakable,colback=teal!4,colframe=teal!35!black,title={Before you move on}]
What does \bn{বিদায়} mean? (\textquotedblleft{}A farewell; a parting\textquotedblright{}.) Say it once more.
\end{tcolorbox}
\section[\texorpdfstring{\bn{শোভাযাত্রা} (shobhājātrā)}{শোভাযাত্রা (shobhājātrā)}]{\bn{শোভাযাত্রা} (shobhājātrā) — a festive procession}

\label{lesson:BN-C308-shobhajatra}



\begin{quote}
Before the new one: say the Bengali for to set afloat, then the Bengali for a farewell.
\end{quote}

\subsection*{You'll want to know: \bn{শোভাযাত্রা}}
\textbf{\bn{শোভাযাত্রা}} — \emph{shobhājātrā} — \textquotedblleft{}a festive procession\textquotedblright{}.

\begin{quote}
\bn{বিসর্জনের} \bn{শোভাযাত্রা} \bn{ঢাক} \bn{বাজিয়ে} \bn{নদীর} \bn{দিকে} \bn{গেল।}
\end{quote}

\begin{quote}
\emph{bisôrjôner shobhājātrā ḍhāk bājiye nôdīr dike gelô.}
\end{quote}

\begin{quote}
The immersion procession went down to the river with the drums beating.
\end{quote}

\begin{cousinweb}
\textbf{\bn{শোভাযাত্রা}} (\emph{shobhājātrā}) joins Sanskrit \emph{śobhā}, 'splendour, beauty' (Monier-Williams), to \textbf{\bn{যাত্রা}}, a journey: a march in finery. A protest march is a \textbf{\bn{মিছিল}}.
\end{cousinweb}

\begin{grammarlens}[title={What you've built}]
One more word for telling what happened.
\end{grammarlens}

\subsection*{Your turn}
\begin{itemize}
\raggedright
  \item \emph{Say it:} \emph{shobhājātrā}
  \item \emph{Say it:} \emph{shobhājātrā}, once more
  \item \emph{Say it:} \emph{bisôrjôner shobhājātrā ḍhāk bājiye nôdīr dike gelô.} — The immersion procession went down to the river with the drums beating.
  \item \emph{Recall:} say the Bengali for to set afloat, then the Bengali for a farewell, then say \emph{shobhājātrā} again
\end{itemize}

\begin{tcolorbox}[breakable,colback=teal!4,colframe=teal!35!black,title={Before you move on}]
What does \bn{শোভাযাত্রা} mean? (\textquotedblleft{}A festive procession\textquotedblright{}.) Say it once more.
\end{tcolorbox}
\section[\texorpdfstring{\bn{কোলাকুলি} \bn{করা} (kolākuli kôrā)}{কোলাকুলি করা (kolākuli kôrā)}]{\bn{কোলাকুলি} \bn{করা} (kolākuli kôrā) — to embrace (in festive greeting)}

\label{lesson:BN-C308-kolakuli-kora}



\begin{quote}
Before the new one: say the Bengali for a farewell, then the Bengali for a festive procession.
\end{quote}

\subsection*{You'll want to know: \bn{কোলাকুলি} \bn{করা}}
\textbf{\bn{কোলাকুলি} \bn{করা}} — \emph{kolākuli kôrā} — \textquotedblleft{}to embrace (in festive greeting)\textquotedblright{}.

\begin{quote}
\bn{বিসর্জনের} \bn{পরে} \bn{বন্ধুরা} \bn{কোলাকুলি} \bn{করল।}
\end{quote}

\begin{quote}
\emph{bisôrjôner pôre bôndhurā kolākuli kôrlô.}
\end{quote}

\begin{quote}
After the immersion the friends embraced one another.
\end{quote}

Past (he/she …): \bn{তিনি} \bn{বন্ধুর} \bn{সঙ্গে} \bn{কোলাকুলি} \bn{করলেন}, \emph{tini bôndhur sôngge kolākuli kôrlen}.

\begin{cousinweb}
\textbf{\bn{কোলাকুলি} \bn{করা}} (\emph{kolākuli kôrā}) doubles \bn{কোল}, the lap or the bosom — CDIAL 3607, Sanskrit \emph{krōḍa}, 'breast', beside \emph{kōla}, 'breast, lap' — into the word for the embrace of the festival's last evening, with \bn{করা}.
\end{cousinweb}

\begin{grammarlens}[title={What you've built}]
One more word for telling what happened.
\end{grammarlens}

\subsection*{Your turn}
\begin{itemize}
\raggedright
  \item \emph{Say it:} \emph{kolākuli kôrā}
  \item \emph{Say it:} \emph{kolākuli kôrā}, once more
  \item \emph{Say it:} \emph{bisôrjôner pôre bôndhurā kolākuli kôrlô.} — After the immersion the friends embraced one another.
  \item \emph{Recall:} say the Bengali for a farewell, then the Bengali for a festive procession, then say \emph{kolākuli kôrā} again
\end{itemize}

\begin{tcolorbox}[breakable,colback=teal!4,colframe=teal!35!black,title={Before you move on}]
What does \bn{কোলাকুলি} \bn{করা} mean? (\textquotedblleft{}To embrace (in festive greeting)\textquotedblright{}.) Say it once more.
\end{tcolorbox}
